(1) 在回旋半径远小于磁场变化尺度时，磁矩 $\mu_{\rm ad}=mv_\perp^2/(2B)$ 是第一绝热不变量。纯磁力与粒子速度垂直，不改变动能，因此 $m(v_\perp^2+v_\parallel^2)/2=mv_0^2/2$ 也保持常量，两式共同给出平行与垂直分速度的能量交换。
(2) 由初态 $v_{\perp0}=v_0\sin\alpha_0$ 及磁矩守恒，$v_\perp^2(B)=v_0^2\sin^2\alpha_0 B/B_0$。再由总速度平方恒定，$v_\parallel^2(B)=v_0^2[1-(B/B_0)\sin^2\alpha_0]$；物理轨道要求后一表达式非负。
(3) 在镜点平行速度趋零，故 $1-(B_m/B_0)\sin^2\alpha_0=0$，解得 $B_m=B_0/\sin^2\alpha_0$。若场沿程最高为 $B_{\max}$，只有满足 $B_m\le B_{\max}$ 的粒子才被反射，即 $\sin^2\alpha_0\ge B_0/B_{\max}$。
(4) 定义磁镜比 $\mathcal R=B_{\max}/B_0$，损失锥边界由 $\alpha_{\rm lc}=\arcsin(1/\sqrt{\mathcal R})$ 给出。初始角小于损失锥半角的粒子无法在该镜中停止平行运动；当 $\mathcal R\to1$，边界趋 $\pi/2$，无增强磁场时不存在普通磁镜反射。
